\documentclass[11pt]{article}
\usepackage[T1]{fontenc}
\usepackage{lmodern}
\usepackage[margin=1in]{geometry}
\usepackage{amsmath,amssymb,amsthm}
\usepackage{booktabs,array}
\usepackage[hidelinks]{hyperref}

\newcommand{\llbracket}{[\![}
\newcommand{\rrbracket}{]\!]}

\hypersetup{
  pdftitle={Shifted Chebyshev Connection Constants and Peakless Motzkin Paths: Proof and Lean 4 Formalization of Ramsey's OEIS A191579 Conjectures},
  pdfauthor={Derek Blair}
}

\newtheorem{theorem}{Theorem}[section]
\newtheorem{proposition}[theorem]{Proposition}
\newtheorem{lemma}[theorem]{Lemma}
\newtheorem{corollary}[theorem]{Corollary}
\theoremstyle{definition}
\newtheorem{definition}[theorem]{Definition}
\newtheorem{remark}[theorem]{Remark}

\newcommand{\N}{\mathbb{N}}
\newcommand{\Z}{\mathbb{Z}}
\newcommand{\seq}[1]{\texttt{#1}}
\newcommand{\StepU}{\mathsf{U}}
\newcommand{\StepH}{\mathsf{H}}
\newcommand{\StepD}{\mathsf{D}}
\newcommand{\Exc}{\mathrm{Exc}}
\newcommand{\Pref}{\mathrm{Pref}}
\newcommand{\PrefH}{\mathrm{PrefH}}

\setlength{\emergencystretch}{2em}

\title{Shifted Chebyshev Connection Constants and Peakless Motzkin Paths:\\[2pt]
\large Proof and Lean~4 Formalization of Ramsey's OEIS \seq{A191579} Conjectures}
\author{Derek Blair}
\date{24 November 2026}

\begin{document}
\maketitle

\begin{abstract}
Let $S(m, j) = \seq{A192062}(m, j)$ ($m, j \ge 0$) be Kenneth J.\ Ramsey's square array defined columnwise by $S(0, j) = 0$, $S(1, j) = 1$, $S(2i, j) = S(2i - 2, j) + S(2i - 1, j)$, and $S(2i + 1, j) = S(2i - 1, j) + j\,S(2i, j)$ ($i \ge 1$), whose even-indexed columns record denominators of continued fraction convergents to $\sqrt{j^2 + 2j}$. Let $R(n, k) = \seq{A191579}(n, k)$ ($0 \le k \le n$, zero-extended for $k > n$) be the $0$-indexed connection-coefficient triangle implicitly determined by the linear relations $S(2n + 2, j + 1) = \sum_{k=0}^{n} R(n, k)\,S(2k + 2, j)$ for all $j \ge 0$. On June~7, 2011, Ramsey conjectured in OEIS \seq{A191579} that (1) the row sums equal \seq{A091964} (left factors of peakless Motzkin paths), (2) the first column equals \seq{A004148} (peakless Motzkin excursions / RNA secondary structure numbers), and (3)--(4) the second and third columns equal \seq{A089735} and \seq{A098075} (the $2$-fold and $3$-fold self-convolutions of \seq{A004148}). R.~J.\ Mathar and P.\ Del\'eham further asked whether the entire triangle coincides with the Riordan array \seq{A097724}, and P.\ Bala (June~26, 2025) conjectured on \seq{A097724} that its entries are the connection constants expanding the shifted rescaled Chebyshev polynomial $u_n(X + 1)$ in the basis $\{u_k(X)\}_{0 \le k \le n}$, where $u_n(X) = U_n(X/2)$.

We resolve all of these claims unconditionally for all $n, k \in \N$. Eliminating the odd rows of $S$ identifies $S(2n + 2, j) = u_n(j + 2)$, where $u_n \in \Z[X]$ is monic of degree $n$, establishing uniqueness of each row $R(n, \cdot)$. Combinatorial last-visit-to-height-$0$ bijections on peakless nonnegative Motzkin words yield the subtraction-free quadratic identity $G(z) + z^2 G(z) = 1 + z G(z) + z^2 G(z)^2$ in $\N\llbracket z \rrbracket$ for $G(z) = \sum_{n \ge 0} \seq{A004148}(n)\,z^n$, together with $\seq{A097724}(n, k) = [z^n](z^k G(z)^{k+1})$. Working purely inside $(\Z[X])\llbracket z \rrbracket$---without formal square roots, power-series substitution, or Lagrange inversion---we factor $1 - X z G + z^2 G^2 = (1 - (X+1)z + z^2)G$ and cancel units to prove $u_n(X + 1) = \sum_{k=0}^n \seq{A097724}(n, k)\,u_k(X)$ in $\Z[X]$, whence $R(n, k) = \seq{A097724}(n, k)$ for all $n, k \ge 0$ and all four \seq{A191579} conjectures follow immediately. The complete proof is formalized and machine-checked in Lean~4 (\texttt{v4.34.1}, Mathlib \texttt{d13f23b7}) with zero \texttt{sorry}s, allowlist-only axioms, and no computational trust escapes.
\end{abstract}

\section{Introduction, Primary Sources, and Prior-Art Boundary}

\subsection{The Continued-Fraction Array \seq{A192062} and Connection Triangle \seq{A191579}}

In June 2011, Kenneth J.\ Ramsey introduced two linked entries in the On-Line Encyclopedia of Integer Sequences: the square array \seq{A192062}~\cite{oeis192062} and the triangular array \seq{A191579}~\cite{oeis191579}. For each column parameter $j \in \N$, the sequence $m \mapsto S(m, j) := \seq{A192062}(m, j)$ is generated by the parity-alternating linear recurrence
\begin{equation}
\label{eq:S-rec}
\begin{aligned}
S(0, j) &= 0, & S(1, j) &= 1, \\
S(2i, j) &= S(2i - 2, j) + S(2i - 1, j) && (i \ge 1), \\
S(2i + 1, j) &= S(2i - 1, j) + j\,S(2i, j) && (i \ge 1).
\end{aligned}
\end{equation}
For $j = 1$, column $S(\cdot, 1)$ is the Fibonacci sequence; for even $j = 2r \ge 2$, the column $S(\cdot, j)$ gives the denominators of the continued fraction convergents to $\sqrt{(r+1)^2 - 1} = \sqrt{r^2 + 2r}$ (whence Ramsey's title ``continued fractions of square root of $(N^2 - 1)$ for $N > 1$'').

Ramsey defined the triangular array \seq{A191579} implicitly by expressing each even row $S(2n+2, j+1)$ in column $j+1$ as a linear combination of the even rows $S(2k+2, j)$ ($0 \le k \le n$) in the preceding column $j$:
\begin{equation}
\label{eq:R-spec}
S(2n + 2, j + 1) = \sum_{k=0}^{n} R(n, k)\,S(2k + 2, j) \qquad (\forall\, j \in \N),
\end{equation}
together with the normalization $R(n, k) = 0$ for $k > n$. (We use $0$-based indices $n = n_{\mathrm{OEIS}} - 1 \ge 0$ and $k = k_{\mathrm{OEIS}} - 1 \in \{0, \dots, n\}$; see Section~\ref{sec:conventions}.) Solving~\eqref{eq:R-spec} row by row yields the triangle shown in Table~\ref{tab:small-n}.

\begin{table}[ht]
\centering
\caption{Rows $n = 0, \dots, 6$ of $R(n, k) = \seq{A191579}(n, k) = \seq{A097724}(n, k)$, alongside row sums \seq{A091964} and the first three columns \seq{A004148}, \seq{A089735}, \seq{A098075}.}
\label{tab:small-n}
\smallskip
{\small
\begin{tabular}{@{}r l r r r r@{}}
\toprule
$n$ & Row $R(n, 0..n)$ & $\sum_{k=0}^n R(n,k)$ & $R(n,0)$ & $R(n,1)$ & $R(n,2)$ \\
    &                  & $= \seq{A091964}(n)$  & $= \seq{A004148}(n)$ & $= \seq{A089735}(n-1)$ & $= \seq{A098075}(n-2)$ \\
\midrule
$0$ & $[1]$ & $1$ & $1$ & --- & --- \\
$1$ & $[1, 1]$ & $2$ & $1$ & $1$ & --- \\
$2$ & $[1, 2, 1]$ & $4$ & $1$ & $2$ & $1$ \\
$3$ & $[2, 3, 3, 1]$ & $9$ & $2$ & $3$ & $3$ \\
$4$ & $[4, 6, 6, 4, 1]$ & $21$ & $4$ & $6$ & $6$ \\
$5$ & $[8, 13, 13, 10, 5, 1]$ & $50$ & $8$ & $13$ & $13$ \\
$6$ & $[17, 28, 30, 24, 15, 6, 1]$ & $121$ & $17$ & $28$ & $30$ \\
\bottomrule
\end{tabular}
}
\end{table}

In the comments of OEIS \seq{A191579} (live revision \#29, October~13, 2025~\cite{oeis191579}), Ramsey observed empirically:
\begin{quote}\itshape
The row sums of this triangle seems to be \seq{A091964} (verified to 12 terms), cf.\ diagonal sums of the triangle \seq{A124428}. The 1st column seems to be \seq{A004148}. The 2nd and 3rd column seems to be \seq{A089735}, and \seq{A098075} (verified to 10 terms).
\end{quote}
In addition, R.~J.\ Mathar (August~1, 2011, on \seq{A191579}) and Philippe Del\'eham (January~26, 2014, on \seq{A097724}) noted that the full triangle \seq{A191579} appears identical (up to offsets) to Emeric Deutsch's peakless-Motzkin triangle \seq{A097724}~\cite{oeis097724}. More recently, Peter Bala posted the following conjecture on \seq{A097724} (June~26, 2025~\cite{oeis097724}):
\begin{quote}\itshape
Conjecture: the row reverse triangle is the triangle of connection constants for expressing the polynomial $u(n,x+1)$ as a linear combination of the polynomials $u(k,x)$, $0 \le k \le n$, where $u(n,x) = U(n,x/2)$ with $U(n,x)$ the $n$-th Chebyshev polynomial of the second kind.
\end{quote}
(Here ``row reverse'' reflects reading the expansion in descending index order $u_n(x+1) = \sum_{j=0}^n c_j u_{n-j}(x)$ with $c_j = \seq{A097724}(n, n-j)$, which is algebraically identical to $u_n(x+1) = \sum_{k=0}^n \seq{A097724}(n, k)\,u_k(x)$.)

\subsection{Convention Audit and Formalization Pitfalls}
\label{sec:conventions}

A careful audit of the seven live OEIS records reveals six concrete convention hazards that must be handled explicitly in any rigorous statement and formalization:
\begin{enumerate}
\item \textbf{Indexing offset ($1$-based vs.\ $0$-based):} \seq{A191579} carries OEIS offset \texttt{1,5} and refers to rows $n_{\mathrm{OEIS}} \ge 1$ and columns $1 \le k_{\mathrm{OEIS}} \le n_{\mathrm{OEIS}}$, whereas \seq{A091964}, \seq{A004148}, \seq{A089735}, \seq{A098075}, and \seq{A097724} all have offset $0$. Setting $n = n_{\mathrm{OEIS}} - 1 \ge 0$ and $k = k_{\mathrm{OEIS}} - 1 \in \{0, \dots, n\}$ aligns all six sequences on a uniform $0$-based grid, with row $0$ equal to $[1]$ and row $3$ equal to $[2, 3, 3, 1]$.
\item \textbf{Argument order in \seq{A191579} vs.\ \seq{A192062}:} The \texttt{\%F} line of \seq{A191579} (June~7, 2011) writes $T(i, 2n)$ with column first, whereas the examples in \seq{A191579} (\texttt{\%e}, writing $T(8, 2) = 56, \dots, T(8, 5) = 329$) and the defining recurrence of \seq{A192062} (June~21 / July~5, 2011) write $T(\text{row}, \text{col})$. We adopt the canonical $(\text{row}, \text{col})$ convention $S(m, j)$ of \seq{A192062}.
\item \textbf{Column-shift guard ($j + 1$ vs.\ $j$):} In Lean's natural numbers $\N$, writing $S(2n+2, j) = \sum_{k=0}^n R(n,k) S(2k+2, j - 1)$ unguarded over all $j \in \N$ truncates $0 - 1$ to $0$, forcing at $(n, j) = (1, 0)$ the false equality $S(4, 0) = R(1, 0) S(2, 0) + R(1, 1) S(4, 0)$, i.e., $2 = 1 \cdot 1 + 1 \cdot 2 = 3$. Matching Ramsey's explicit restriction $j_{\mathrm{orig}} > 0$, we state~\eqref{eq:R-spec} with $S(2n+2, j+1)$ on the left and $S(2k+2, j)$ on the right for all $j \in \N$.
\item \textbf{Exact domain guards for shifted columns:} In $0$-based coordinates, row $n$ is supported on $0 \le k \le n$ and zero-extended for $k > n$. Thus $R(0, 1) = 0 \ne 1 = \seq{A089735}(0)$ and $R(0, 2) = R(1, 2) = 0 \ne 1 = \seq{A098075}(0)$. The domain guards $n \ge 1$ for $R(n, 1) = \seq{A089735}(n - 1)$ and $n \ge 2$ for $R(n, 2) = \seq{A098075}(n - 2)$ are therefore exact.
\item \textbf{Auxiliary OEIS formula errata:} On \seq{A089735}~\cite{oeis089735}, one contributed \texttt{\%F} line prints \texttt{G.f.: z\textasciicircum 3*S\textasciicircum 2} with an extraneous $z^3$ factor contradicting its own offset $0$ and initial terms $1, 2, 3, 6, 13, \dots$; we define \seq{A089735} directly from its canonical \texttt{\%N} specification as the $2$-fold Cauchy convolution of \seq{A004148}. On \seq{A004148}~\cite{oeis004148}, the \texttt{\%N} recurrence $a(n+1) = a(n) + \sum_{k=1}^{n-1} a(k) a(n-1-k)$ starts at $k = 1$ (not $k = 0$); we define \seq{A004148} combinatorially as peakless Motzkin excursions and derive the recurrence as Proposition~\ref{prop:a004148-rec}.
\item \textbf{Equivalence of left factors and nonnegative prefixes for \seq{A091964}:} Any prefix of a peakless Motzkin excursion stays weakly above the $x$-axis and contains no $\StepU\StepD$ peak. Conversely, any peakless nonnegative $(\StepU, \StepH, \StepD)$-word $w$ ending at height $k \ge 0$ extends to a peakless excursion via $w$ (if $k = 0$) or $w\,\StepH\,\StepD^k$ (if $k \ge 1$, where the inserted $\StepH$ prevents any $\StepU\StepD$ junction). Thus the \texttt{\%N} and \texttt{\%C} definitions of \seq{A091964}~\cite{oeis091964} coincide identically for all $n \in \N$.
\end{enumerate}

\subsection{Prior-Art Boundary}

Literature audits conducted on 14~November~2026 (25 logged queries across OEIS, arXiv, OpenAlex, and web search, including full-text TeX/XML inspection) and re-verified live on 24~November~2026 (live Google web search, OEIS revision check, arXiv API, and OpenAlex) establish the exact novelty boundary:
\begin{itemize}
\item \textbf{Established on the lattice-path / Riordan side alone:} The peakless Motzkin excursion sequence \seq{A004148} (RNA secondary structure numbers) was analyzed by Stein and Waterman~\cite{stein1979}, Hofacker, Schuster, and Stadler~\cite{hofacker1998}, and in lattice-path / kernel-method treatments by Asinowski et al.~\cite{asinowski2019} and Prodinger~\cite{prodinger2023,prodinger2025}. On the Riordan-group side, the triangle \seq{A097724} $= (G(z), z G(z))$, its row sums \seq{A091964} with generating function $G/(1 - z G)$, its columns $G, z G^2, z^2 G^3$, and its pseudo-involution structure appear in Deutsch~\cite{oeis097724}, Cameron and Nkwanta~\cite{cameron2005}, Barry~\cite{barry2022}, and Tenn~\cite{tenn2026}.
\item \textbf{Unproved prior to this work:} Neither \seq{A191579} nor \seq{A192062} appears in any paper indexed by arXiv, OpenAlex, or Google web search (including full-text inspection of~\cite{barry2022,prodinger2025,tenn2026}). No published proof connects Ramsey's continued-fraction array \seq{A192062} or connection-coefficient triangle \seq{A191579} to \seq{A097724} or proves any of the four conjectural identifications in \seq{A191579} (\seq{A091964}, \seq{A004148}, \seq{A089735}, \seq{A098075}); moreover, Bala's June~26, 2025 shifted-Chebyshev connection-constant characterization on \seq{A097724} remains explicitly labeled as a conjecture on the live OEIS record as of 24~November~2026.
\end{itemize}

\subsection{Main Results}

\begin{theorem}[Main Theorem]
\label{thm:main}
Let $S(m, j) = \seq{A192062}(m, j)$ be defined by~\eqref{eq:S-rec}, and let $u_n(X) \in \Z[X]$ be defined by $u_0(X) = 1$, $u_1(X) = X$, and $u_{n+2}(X) = X u_{n+1}(X) - u_n(X)$ for $n \ge 0$. Then:
\begin{enumerate}
\item[\textup{(0)}] \textbf{Existence and uniqueness:} For every $n \in \N$, there exists a unique zero-extended sequence $r : \N \to \N$ (with $r(k) = 0$ for all $k > n$) satisfying~\eqref{eq:R-spec}. Denoting this unique row by $R(n, \cdot) = \seq{A191579}(n, \cdot)$, we have
\[
R(n, k) = \seq{A097724}(n, k) = [z^n]\bigl(z^k G(z)^{k+1}\bigr) \qquad (\forall\, n, k \in \N),
\]
where $G(z) = \sum_{m \ge 0} \seq{A004148}(m)\,z^m \in \N\llbracket z \rrbracket$, and in $\Z[X]$,
\[
u_n(X + 1) = \sum_{k=0}^{n} R(n, k)\,u_k(X) \qquad (\forall\, n \in \N).
\]
\item[\textup{(1)}] \textbf{Row sums:} $\sum_{k=0}^{n} R(n, k) = \seq{A091964}(n)$ for every $n \in \N$.
\item[\textup{(2)}] \textbf{First column ($k = 0$):} $R(n, 0) = \seq{A004148}(n)$ for every $n \in \N$.
\item[\textup{(3)}] \textbf{Second column ($k = 1$):} $R(n, 1) = \seq{A089735}(n - 1)$ for every $n \ge 1$.
\item[\textup{(4)}] \textbf{Third column ($k = 2$):} $R(n, 2) = \seq{A098075}(n - 2)$ for every $n \ge 2$.
\end{enumerate}
\end{theorem}

\section{Reduction of \seq{A192062} to Monic Rescaled Chebyshev Polynomials and Row Uniqueness}

\begin{lemma}[Elimination of odd rows of \seq{A192062}]
\label{lem:S-even-rec}
For all $n, j \in \N$, the even rows of $S = \seq{A192062}$ satisfy the three-term recurrence in $\Z$:
\begin{equation}
\label{eq:S-even-rec}
S(2(n+2)+2, j) = (j + 2)\,S(2(n+1)+2, j) - S(2n+2, j).
\end{equation}
\end{lemma}

\begin{proof}
From~\eqref{eq:S-rec}, for any $n, j \in \N$ we have
\begin{align*}
S(2n+6, j) &= S(2n+4, j) + S(2n+5, j), \\
S(2n+5, j) &= S(2n+3, j) + j\,S(2n+4, j), \\
S(2n+4, j) &= S(2n+2, j) + S(2n+3, j).
\end{align*}
Eliminating $S(2n+3, j) = S(2n+4, j) - S(2n+2, j)$ and $S(2n+5, j)$ in $\Z$ gives
\begin{align*}
S(2n+6, j) &= S(2n+4, j) + \bigl(S(2n+4, j) - S(2n+2, j)\bigr) + j\,S(2n+4, j) \\
           &= (j+2)\,S(2n+4, j) - S(2n+2, j),
\end{align*}
which is~\eqref{eq:S-even-rec}.
\end{proof}

Define the monic rescaled Chebyshev polynomials of the second kind $u_n(X) \in \Z[X]$ ($u_n(X) = U_n(X/2)$) by
\begin{equation}
\label{eq:uPoly}
u_0(X) = 1, \qquad u_1(X) = X, \qquad u_{n+2}(X) = X\,u_{n+1}(X) - u_n(X) \quad (n \ge 0).
\end{equation}

\begin{proposition}[Chebyshev evaluation of \seq{A192062}]
\label{prop:S-eq-uPoly}
For all $n, j \in \N$,
\[
S(2n + 2, j) = u_n(j + 2) \qquad \text{in } \Z.
\]
\end{proposition}

\begin{proof}
We proceed by strong induction on $n$. For $n = 0$, $S(2, j) = S(0, j) + S(1, j) = 1 = u_0(j + 2)$. For $n = 1$, $S(3, j) = S(1, j) + j S(2, j) = 1 + j$, hence $S(4, j) = S(2, j) + S(3, j) = j + 2 = u_1(j + 2)$. For $n = m + 2$, Lemma~\ref{lem:S-even-rec} and the induction hypothesis give
\[
S(2(m+2)+2, j) = (j + 2)\,u_{m+1}(j + 2) - u_m(j + 2) = u_{m+2}(j + 2). \qedhere
\]
\end{proof}

\begin{lemma}[Monicity, linear independence, and row uniqueness]
\label{lem:uniqueness}
\begin{enumerate}
\item[\textup{(a)}] For every $n \in \N$, the polynomial $u_n(X) \in \Z[X]$ is monic of exact degree $n$.
\item[\textup{(b)}] If $m \in \N$ and $c_0, \dots, c_{m-1} \in \Z$ satisfy $\sum_{k=0}^{m-1} c_k u_k(X) = 0$ in $\Z[X]$, then $c_k = 0$ for all $0 \le k < m$.
\item[\textup{(c)}] For every $n \in \N$, at most one function $r : \N \to \N$ satisfies $r(k) = 0$ for all $k > n$ together with~\eqref{eq:R-spec}.
\end{enumerate}
\end{lemma}

\begin{proof}
(a) Strong induction on $n$: $u_0 = 1$ (monic, degree $0$) and $u_1 = X$ (monic, degree $1$). For $n = m + 2$, $X u_{m+1}$ is monic of degree $m + 2$, while $\deg u_m = m < m + 2$, so $u_{m+2} = X u_{m+1} - u_m$ is monic of degree $m + 2$.

(b) Induction on $m$. The case $m = 0$ is vacuous. For $m + 1$, extracting the coefficient of $X^m$ from $\sum_{k=0}^{m} c_k u_k(X) = 0$ annihilates all terms $k < m$ (since $\deg u_k = k < m$) and yields $c_m \cdot [X^m]u_m(X) = c_m \cdot 1 = 0$ by (a). Dropping the $k = m$ term and applying the induction hypothesis yields $c_k = 0$ for all $k < m + 1$.

(c) Suppose $r_1, r_2 : \N \to \N$ both satisfy $r_i(k) = 0$ ($k > n$) and~\eqref{eq:R-spec}. Set $P(X) := \sum_{k=0}^n (r_1(k) - r_2(k))\,u_k(X) \in \Z[X]$. By Proposition~\ref{prop:S-eq-uPoly} and~\eqref{eq:R-spec}, for every $j \in \N$ we have
\[
P(j + 2) = \sum_{k=0}^n r_1(k)\,S(2k + 2, j) - \sum_{k=0}^n r_2(k)\,S(2k + 2, j) = S(2n+2, j+1) - S(2n+2, j+1) = 0.
\]
Since $\{j + 2 : j \in \N\} \subseteq \Z$ is infinite and $\Z$ is an integral domain, $P = 0$ in $\Z[X]$. By (b), $r_1(k) = r_2(k)$ for all $0 \le k \le n$, and for $k > n$ both sides vanish by zero-extension. Hence $r_1 = r_2$.
\end{proof}

\section{Combinatorial Decompositions on Peakless Motzkin Words}

Let $\Sigma = \{\StepU, \StepH, \StepD\}$ be the Motzkin step alphabet with vertical displacements $\delta(\StepU) = 1$, $\delta(\StepH) = 0$, $\delta(\StepD) = -1$. For a word $w = w_0 w_1 \cdots w_{L-1} \in \Sigma^*$ and $0 \le i \le L$, let $h(w, i) := \sum_{t=0}^{i-1} \delta(w_t)$ be the height after $i$ steps ($h(w, 0) = 0$). A word $w$ is \emph{peakless} if there is no index $i$ with $i + 1 < |w|$, $w_i = \StepU$, and $w_{i+1} = \StepD$. It is a \emph{peakless nonnegative prefix} if it is peakless and $h(w, i) \ge 0$ for every $0 \le i \le |w|$.

For $n, k \in \N$, define the finite sets
\begin{align*}
\Pref(n) &:= \{w \in \Sigma^n : w \text{ is a peakless nonnegative prefix}\}, \\
\PrefH(n, k) &:= \{w \in \Pref(n) : h(w, n) = k\}, \\
\Exc(n) &:= \PrefH(n, 0),
\end{align*}
so that $\seq{A091964}(n) = |\Pref(n)|$, $\seq{A097724}(n, k) = |\PrefH(n, k)|$, and $\seq{A004148}(n) = |\Exc(n)|$. Furthermore, by definition,
\begin{align*}
\seq{A089735}(n) &= \sum_{k=0}^{n} \seq{A004148}(k)\,\seq{A004148}(n - k), \\
\seq{A098075}(n) &= \sum_{i=0}^{n}\sum_{j=0}^{n-i} \seq{A004148}(i)\,\seq{A004148}(j)\,\seq{A004148}(n - i - j).
\end{align*}

Two elementary boundary constraints govern peakless concatenation:
\begin{itemize}
\item If $w \in \Pref(L)$ with $L > 0$, then $0 \le h(w, 1) = \delta(w_0)$, so $w_0 \ne \StepD$.
\item If $w \in \Exc(L)$ with $L > 0$, then $0 = h(w, L) = h(w, L-1) + \delta(w_{L-1})$ with $h(w, L-1) \ge 0$, so $\delta(w_{L-1}) \le 0$ and hence $w_{L-1} \ne \StepU$.
\end{itemize}
Moreover, for any $w \in \Pref(N+1)$ with $h(w, N+1) > 0$, the set $\{0 \le i \le N+1 : h(w, i) = 0\}$ contains $0$ and excludes $N+1$, so its maximum $m = m(w)$ satisfies $0 \le m \le N$, $h(w, m) = 0$, $h(w, i) \ge 1$ for all $m < i \le N+1$, and consequently $\delta(w_m) = h(w, m+1) - h(w, m) \ge 1$, forcing $w_m = \StepU$.

\begin{proposition}[Excursion decomposition and quadratic generating-function equation]
\label{prop:a004148-rec}
Write $a(n) := \seq{A004148}(n)$. Then $a(0) = a(1) = 1$, and for every $n \in \N$,
\begin{equation}
\label{eq:a004148-rec}
a(n + 2) = a(n + 1) + \sum_{k=0}^{n-1} a(k + 1)\,a(n - 1 - k).
\end{equation}
Consequently, the ordinary generating function $G(z) := \sum_{n \ge 0} a(n)\,z^n \in \N\llbracket z \rrbracket$ satisfies the subtraction-free identity
\begin{equation}
\label{eq:G-quad}
G(z) + z^2 G(z) = 1 + z G(z) + z^2 G(z)^2.
\end{equation}
\end{proposition}

\begin{proof}
Clearly $\Exc(0) = \{\varepsilon\}$ and $\Exc(1) = \{\StepH\}$, giving $a(0) = a(1) = 1$. For any $n \in \N$, we construct a bijection
\[
\Exc(n + 2) \;\xrightarrow{\;\sim\;}\; \Exc(n + 1) \;\sqcup\; \bigsqcup_{k=0}^{n-1} \Exc(k + 1) \times \Exc(n - 1 - k).
\]
Given $w \in \Exc(n+2)$, its final letter satisfies $w_{n+1} \ne \StepU$ by the boundary constraint above.
\begin{itemize}
\item If $w_{n+1} = \StepH$, write $w = w' \StepH$; then $w' \in \Exc(n + 1)$, and conversely appending $\StepH$ to any $w' \in \Exc(n + 1)$ preserves nonnegativity, peaklessness ($\StepH \ne \StepD$), and endpoint height $0$.
\item If $w_{n+1} = \StepD$, then $h(w, n+1) = 1 > 0$. Let $m = \max\{0 \le i \le n+1 : h(w, i) = 0\} \le n$. As noted above, $w_m = \StepU$ and $h(w, i) \ge 1$ for all $m < i \le n+1$. Since $w$ is peakless and $w_{n+1} = \StepD$, we cannot have $m = n$; thus $m \le n - 1$. Setting $k := n - 1 - m \in \{0, \dots, n-1\}$ decomposes $w$ uniquely as
\[
w = w_1 \,\StepU\, w_2 \,\StepD,
\]
where $w_1$ has length $m = n - 1 - k$ and $w_2$ has length $n - m = k + 1 \ge 1$. For $0 \le r \le k + 1$, $h(w, m + 1 + r) = 1 + h(w_2, r) \ge 1$, and at $r = k + 1$ we obtain $1 + h(w_2, k+1) = h(w, n+1) = 1$; hence $w_1 \in \Exc(n - 1 - k)$ and $w_2 \in \Exc(k + 1)$.

Conversely, given $0 \le k < n$, $w_2 \in \Exc(k + 1)$, and $w_1 \in \Exc(n - 1 - k)$, the concatenation $w = w_1\,\StepU\,w_2\,\StepD$ has length $n + 2$, ends in $\StepD$, has height $0$ at index $m_0 := |w_1| = n - 1 - k$, and has height $1 + h(w_2, r) \ge 1$ at every index $m_0 < i \le n + 1$; therefore $m_0$ is uniquely recovered as $\max\{0 \le i \le n+1 : h(w, i) = 0\}$. Peaklessness holds across all three junctions because $\StepU \ne \StepD$, $w_2$ is nonempty with first letter $\ne \StepD$, and $w_2$ is an excursion with last letter $\ne \StepU$.
\end{itemize}
Taking cardinalities yields~\eqref{eq:a004148-rec}.

For~\eqref{eq:G-quad}, we compare coefficients of $z^m$ in $\N$:
\begin{itemize}
\item At $m = 0$: $[z^0](G + z^2 G) = a(0) = 1 = [z^0](1 + z G + z^2 G^2)$.
\item At $m = 1$: $[z^1](G + z^2 G) = a(1) = 1 = a(0) = [z^1](1 + z G + z^2 G^2)$.
\item At $m = n + 2$ ($n \ge 0$): splitting off the $j = 0$ term from $[z^n](G^2) = \sum_{j=0}^n a(j) a(n - j)$ and using $a(0) = 1$ gives
\begin{align*}
[z^{n+2}](1 + z G + z^2 G^2)
  &= a(n+1) + a(n) + \sum_{k=0}^{n-1} a(k+1)\,a(n-1-k) \\
  &= a(n+2) + a(n) = [z^{n+2}](G + z^2 G),
\end{align*}
by~\eqref{eq:a004148-rec}. \qedhere
\end{itemize}
\end{proof}

\begin{proposition}[Prefix decomposition and Riordan columns of \seq{A097724}]
\label{prop:a097724-gf}
Write $P(n, k) := \seq{A097724}(n, k)$. For all $n, k \in \N$,
\begin{equation}
\label{eq:a097724-gf}
P(n, k) = [z^n]\bigl(z^k G(z)^{k+1}\bigr),
\end{equation}
and hence $P(n, k) = 0$ whenever $k > n$. Moreover, for every $n \in \N$,
\begin{equation}
\label{eq:a091964-sum}
\seq{A091964}(n) = \sum_{k=0}^{n} P(n, k).
\end{equation}
\end{proposition}

\begin{proof}
We first establish the column recurrence
\begin{equation}
\label{eq:P-rec}
P(0, k + 1) = 0, \qquad P(n + 1, k + 1) = \sum_{m=0}^{n} a(m)\,P(n - m, k) \qquad (\forall\, n, k \in \N).
\end{equation}
The equality $P(0, k + 1) = 0$ is immediate because the only word of length $0$ has endpoint height $0 \ne k + 1$. For any $n, k \in \N$, we construct a bijection
\[
\PrefH(n + 1, k + 1) \;\xrightarrow{\;\sim\;}\; \bigsqcup_{m=0}^{n} \Exc(m) \times \PrefH(n - m, k).
\]
Given $w \in \PrefH(n + 1, k + 1)$, since $h(w, n+1) = k + 1 > 0$, let $m = \max\{0 \le i \le n+1 : h(w, i) = 0\} \le n$. Then $w_m = \StepU$ and $h(w, i) \ge 1$ for $m < i \le n + 1$, splitting $w = w_1\,\StepU\,w_2$ with $w_1 \in \Exc(m)$ and $w_2 \in \PrefH(n - m, k)$ (since $h(w, m+1+r) = 1 + h(w_2, r) \ge 1$ and $1 + h(w_2, n-m) = k+1$). Conversely, for any $0 \le m \le n$, $w_1 \in \Exc(m)$, and $w_2 \in \PrefH(n - m, k)$, the word $w = w_1\,\StepU\,w_2$ has length $n + 1$, endpoint height $k + 1$, no peak across either junction (as $\StepU \ne \StepD$, and if $w_2$ is nonempty then its first letter is $\ne \StepD$), and unique last height-$0$ index $m$. Taking cardinalities gives~\eqref{eq:P-rec}.

Let $P_k(z) := \sum_{n \ge 0} P(n, k)\,z^n \in \N\llbracket z \rrbracket$. Then $P_0(z) = G(z)$ since $\PrefH(n, 0) = \Exc(n)$, and~\eqref{eq:P-rec} asserts precisely that $P_{k+1}(z) = z G(z) P_k(z)$ for all $k \in \N$. Induction on $k$ yields $P_k(z) = z^k G(z)^{k+1}$, which is~\eqref{eq:a097724-gf}.

Finally, every $w \in \Pref(n)$ satisfies $0 \le h(w, n) \le n$ (since $\delta(x) \le 1$ for every step $x \in \Sigma$), partitioning $\Pref(n) = \bigsqcup_{k=0}^n \PrefH(n, k)$ by endpoint height and proving~\eqref{eq:a091964-sum}.
\end{proof}

\section{Square-Root-Free Generating-Function Proof in \texorpdfstring{$(\Z[X])\llbracket z \rrbracket$}{(Z[X])[[z]]}}

We now lift $G(z) \in \N\llbracket z \rrbracket$ via the canonical ring homomorphism $\N \to \Z[X]$ to a power series $G(z) \in (\Z[X])\llbracket z \rrbracket$, which still satisfies~\eqref{eq:G-quad} and $[z^n](z^k G(z)^{k+1}) = P(n, k) \in \Z[X]$ (as constant polynomials). Define two formal power series in $(\Z[X])\llbracket z \rrbracket$:
\begin{equation}
\label{eq:VW-def}
V(z) := \sum_{n=0}^{\infty} u_n(X + 1)\,z^n,
\qquad
W(z) := \sum_{n=0}^{\infty} \left(\sum_{k=0}^{n} P(n, k)\,u_k(X)\right) z^n.
\end{equation}

\begin{lemma}[Rational generating function for $V(z)$]
\label{lem:V-denom}
In $(\Z[X])\llbracket z \rrbracket$,
\[
\bigl(1 - (X + 1)z + z^2\bigr)\,V(z) = 1.
\]
\end{lemma}

\begin{proof}
Comparing coefficients of $z^m$ in $V(z) - (X+1) z V(z) + z^2 V(z)$:
\begin{itemize}
\item $m = 0$: $u_0(X+1) = 1$.
\item $m = 1$: $u_1(X+1) - (X+1) u_0(X+1) = (X+1) - (X+1) = 0$.
\item $m = n + 2$: $u_{n+2}(X+1) - (X+1) u_{n+1}(X+1) + u_n(X+1) = 0$, obtained by composing~\eqref{eq:uPoly} with $X + 1$. \qedhere
\end{itemize}
\end{proof}

\begin{lemma}[Shift-power action on $W(z)$]
\label{lem:shift-W}
For all $d, n \in \N$,
\[
[z^n]\bigl((z G(z))^d\, W(z)\bigr) = \sum_{k=0}^{n} P(n, k + d)\,u_k(X).
\]
\end{lemma}

\begin{proof}
By definition of Cauchy multiplication,
\[
[z^n]\bigl((z G)^d W\bigr) = \sum_{i + j = n} [z^i]((z G)^d) \cdot [z^j]W.
\]
For every summand $i + j = n$ we have $j \le n$, and since $P(j, k) = 0$ for $k > j$, we may extend the finite sum $[z^j]W = \sum_{k=0}^j P(j, k) u_k(X)$ to $0 \le k \le n$ and replace $P(j, k)$ by $[z^j](z^k G^{k+1})$ via Proposition~\ref{prop:a097724-gf}. Exchanging the two finite sums gives
\begin{align*}
[z^n]\bigl((z G)^d W\bigr)
&= \sum_{k=0}^{n} \left(\sum_{i + j = n} [z^i]((z G)^d) \cdot [z^j](z^k G^{k+1})\right) u_k(X) \\
&= \sum_{k=0}^{n} [z^n]\bigl((z G)^d \cdot z^k G^{k+1}\bigr)\,u_k(X) \\
&= \sum_{k=0}^{n} [z^n]\bigl(z^{k+d} G^{k+d+1}\bigr)\,u_k(X)
 = \sum_{k=0}^{n} P(n, k + d)\,u_k(X). \qedhere
\end{align*}
\end{proof}

\begin{lemma}[Quadratic annihilator of $W(z)$]
\label{lem:W-quad}
In $(\Z[X])\llbracket z \rrbracket$,
\begin{equation}
\label{eq:W-quad}
\bigl(1 - X z G(z) + z^2 G(z)^2\bigr)\,W(z) = G(z).
\end{equation}
\end{lemma}

\begin{proof}
By Lemma~\ref{lem:shift-W} applied with $d = 0, 1, 2$, for any $n \in \N$ the $z^n$-coefficient of the left-hand side of~\eqref{eq:W-quad} is
\[
\sum_{k=0}^{n} P(n, k)\,u_k(X) - X \sum_{k=0}^{n} P(n, k+1)\,u_k(X) + \sum_{k=0}^{n} P(n, k+2)\,u_k(X).
\]
Since $P(n, j) = 0$ for all $j > n$:
\begin{itemize}
\item For $n = 0$: the expression equals $P(0, 0) u_0(X) = P(0, 0) = a(0)$.
\item For $n = 1$: it equals $P(1, 0) u_0(X) + P(1, 1) (u_1(X) - X u_0(X)) = P(1, 0) = a(1)$.
\item For $n = m + 2$: dropping the vanishing top terms $P(m+2, m+3) = P(m+2, m+4) = 0$ and re-indexing the three sums leaves
\begin{align*}
&P(m+2, 0)\,u_0(X) + P(m+2, 1)\bigl(u_1(X) - X u_0(X)\bigr) \\
&\qquad + \sum_{k=0}^{m} P(m+2, k+2)\bigl(u_{k+2}(X) - X u_{k+1}(X) + u_k(X)\bigr) \\
&= P(m+2, 0) = a(m+2) = [z^{m+2}] G(z),
\end{align*}
because $u_0(X) = 1$, $u_1(X) - X u_0(X) = 0$, and $u_{k+2}(X) - X u_{k+1}(X) + u_k(X) = 0$ by~\eqref{eq:uPoly}. \qedhere
\end{itemize}
\end{proof}

\begin{theorem}[Shifted Chebyshev expansion and equality $W(z) = V(z)$]
\label{thm:W-eq-V}
In $(\Z[X])\llbracket z \rrbracket$, $W(z) = V(z)$. Consequently, for every $n \in \N$,
\begin{equation}
\label{eq:bala-identity}
u_n(X + 1) = \sum_{k=0}^{n} \seq{A097724}(n, k)\,u_k(X) \qquad \text{in } \Z[X].
\end{equation}
\end{theorem}

\begin{proof}
Rearranging the quadratic identity~\eqref{eq:G-quad} in $(\Z[X])\llbracket z \rrbracket$ gives $1 + z^2 G^2 = G - z G + z^2 G$. Subtracting $X z G$ from both sides yields the exact ring factorization
\begin{equation}
\label{eq:factorization}
1 - X z G(z) + z^2 G(z)^2 = \bigl(1 - (X + 1)z + z^2\bigr)\,G(z).
\end{equation}
Substituting~\eqref{eq:factorization} into Lemma~\ref{lem:W-quad} gives
\[
\bigl(1 - (X + 1)z + z^2\bigr)\,W(z)\,G(z) = 1 \cdot G(z).
\]
In the power series ring $A\llbracket z \rrbracket$ over any commutative ring $A$, a series is a unit if and only if its constant coefficient is a unit in $A$. Since $[z^0]G(z) = a(0) = 1$ and $[z^0](1 - (X+1)z + z^2) = 1$ are units in $\Z[X]$, both $G(z)$ and $1 - (X+1)z + z^2$ are units in $(\Z[X])\llbracket z \rrbracket$. Cancelling $G(z)$ gives $(1 - (X+1)z + z^2)\,W(z) = 1 = (1 - (X+1)z + z^2)\,V(z)$ (by Lemma~\ref{lem:V-denom}); cancelling $1 - (X+1)z + z^2$ gives $W(z) = V(z)$. Extracting $[z^n]$ gives~\eqref{eq:bala-identity}.
\end{proof}

\begin{proof}[Proof of Theorem~\ref{thm:main}]
Fix $n \in \N$. Evaluating~\eqref{eq:bala-identity} at $X = j + 2$ ($j \in \N$) and applying Proposition~\ref{prop:S-eq-uPoly} gives
\[
S(2n + 2, j + 1) = u_n((j + 1) + 2) = u_n((j + 2) + 1) = \sum_{k=0}^{n} P(n, k)\,u_k(j + 2) = \sum_{k=0}^{n} P(n, k)\,S(2k + 2, j).
\]
Since $P(n, k) = 0$ for $k > n$ (Proposition~\ref{prop:a097724-gf}), the row $k \mapsto P(n, k)$ satisfies~\eqref{eq:R-spec} and zero-extension. Combined with Lemma~\ref{lem:uniqueness}(c), this proves existence and uniqueness of row $n$ of $R$, and establishes $R(n, k) = P(n, k) = \seq{A097724}(n, k) = [z^n](z^k G(z)^{k+1})$ for all $n, k \in \N$, proving part~(0) and~\eqref{eq:bala-identity}.

Parts~(1)--(4) now follow directly:
\begin{enumerate}
\item[\textup{(1)}] $\sum_{k=0}^n R(n, k) = \sum_{k=0}^n P(n, k) = \seq{A091964}(n)$ by~\eqref{eq:a091964-sum}.
\item[\textup{(2)}] $R(n, 0) = P(n, 0) = \seq{A004148}(n)$ for all $n \in \N$.
\item[\textup{(3)}] For $n \ge 1$, writing $n = (n - 1) + 1$ gives
\[
R(n, 1) = [z^{(n-1)+1}]\bigl(z G(z)^2\bigr) = [z^{n-1}]\bigl(G(z)^2\bigr) = \sum_{k=0}^{n-1} a(k)\,a(n - 1 - k) = \seq{A089735}(n - 1).
\]
\item[\textup{(4)}] For $n \ge 2$, writing $n = (n - 2) + 2$ gives
\begin{align*}
R(n, 2) &= [z^{(n-2)+2}]\bigl(z^2 G(z)^3\bigr) = [z^{n-2}]\bigl(G(z)\cdot(G(z)\cdot G(z))\bigr) \\
&= \sum_{i=0}^{n-2}\sum_{j=0}^{n-2-i} a(i)\,a(j)\,a(n - 2 - i - j) = \seq{A098075}(n - 2). \qedhere
\end{align*}
\end{enumerate}
\end{proof}

\begin{remark}[Avoidance of formal square roots and Riordan substitution]
The classical Riordan-group proof of a connection-constant identity for $(G(z), z G(z))$ composes power series via $V(z) = G(z) / (1 - X z G(z) + z^2 G(z)^2)$ using formal substitution $t = z G(z)$ or the explicit algebraic branch $G(z) = \bigl(1 - z + z^2 - \sqrt{1 - 2z - z^2 - 2z^3 + z^4}\bigr)/(2z^2)$. By working directly with the bivariate polynomial-coefficient series $W(z)$ and factoring $1 - X z G + z^2 G^2 = (1 - (X+1)z + z^2)G$ via~\eqref{eq:G-quad}, Theorem~\ref{thm:W-eq-V} uses only ring multiplication, finite sum re-indexing, and unit cancellation in $(\Z[X])\llbracket z \rrbracket$.
\end{remark}

\section{Lean 4 Formalization and Verification Record}

The formalization resides in the repository \texttt{derekblair/math-proofs} under Lean~\texttt{v4.34.1} (\nolinkurl{leanprover/lean4:v4.34.1}) and Mathlib revision
{\small\nolinkurl{d13f23b723b8a846827a245b89c10fc7d3f11612}}:
\begin{itemize}
\item \nolinkurl{MathProofs/SeamB/C042/Statement.lean} (119 lines; SHA-256:\\
{\small\nolinkurl{5ee8b6f0acfb217cf7ee550361c575c5fc53429dc47357717b41b4b89a3fd239}}):\\
transcribed byte-for-byte from the frozen manifest \texttt{PIN-C042.md} v1 (\texttt{delta\_vs\_source: NONE}, SHA-256:\\
{\small\nolinkurl{2c68f4c0ab89c832ab06bc973758ff53434148b319516b8ee33a745d2f31c2e4}}).
\item \nolinkurl{MathProofs/SeamB/C042/Solution.lean} (1{,}417 lines; SHA-256:\\
{\small\nolinkurl{d4ed10666f54e2167fecc59caf6d3b4b7bedab0d26bed060ca2d7f044360aaa2}}):\\
contains the verbatim definitions and exact theorem signatures of \texttt{Statement.lean} in \nolinkurl{namespace MathProofs.SeamB.C042}, closing all five pinned theorems plus the full triangle equality \texttt{a191579\_eq\_a097724} with zero \texttt{sorry}s, zero warnings, and zero custom or computational axioms.
\end{itemize}

Table~\ref{tab:lean-map} records the correspondence between the mathematical results in this paper and the checked declarations in \nolinkurl{MathProofs/SeamB/C042/Solution.lean}.

\begin{table}[ht]
\centering
\caption{Correspondence between paper results and Lean~4 declarations in \texttt{MathProofs.SeamB.C042}.}
\label{tab:lean-map}
\smallskip
{\small
\begin{tabular}{@{}l >{\raggedright\arraybackslash}p{8.2cm} l@{}}
\toprule
Paper Result & Lean~4 Declaration(s) & Status \\
\midrule
Lemma~\ref{lem:S-even-rec} & \texttt{a192062\_even\_recurrence} & Private, proved \\
Proposition~\ref{prop:S-eq-uPoly} & \texttt{a192062\_even\_eq\_uPoly} & Private, proved \\
Lemma~\ref{lem:uniqueness} & \texttt{uPoly\_monic\_and\_natDegree}, \texttt{isA191579Row\_unique} & Private, proved \\
Proposition~\ref{prop:a004148-rec} & \texttt{excDecompEquiv}, \texttt{a004148\_recurrence}, \texttt{G\_nat\_quad} & Private, proved \\
Proposition~\ref{prop:a097724-gf} & \texttt{prefHDecompEquiv}, \texttt{a097724\_eq\_coeff\_X\_pow\_mul\_G\_pow} & Private, proved \\
Lemmas~\ref{lem:V-denom}--\ref{lem:W-quad} & \texttt{V\_ser\_mul\_denom}, \texttt{coeff\_ZG\_pow\_mul\_W}, \texttt{W\_ser\_quad\_eq} & Private, proved \\
Theorem~\ref{thm:W-eq-V} & \texttt{W\_ser\_eq\_V\_ser}, \texttt{isA191579Row\_a097724} & Private, proved \\
Theorem~\ref{thm:main}(0) & \texttt{a191579\_row\_exists\_unique} (pinned), \texttt{a191579\_eq\_a097724} & Exported, proved \\
Theorem~\ref{thm:main}(1) & \texttt{a191579\_rowSum\_eq\_a091964} (pinned) & Exported, proved \\
Theorem~\ref{thm:main}(2) & \texttt{a191579\_col0\_eq\_a004148} (pinned) & Exported, proved \\
Theorem~\ref{thm:main}(3) & \texttt{a191579\_col1\_eq\_a089735} (pinned, guard $1 \le n$) & Exported, proved \\
Theorem~\ref{thm:main}(4) & \texttt{a191579\_col2\_eq\_a098075} (pinned, guard $2 \le n$) & Exported, proved \\
\bottomrule
\end{tabular}
}
\end{table}

\noindent\textbf{Audit summary.}
\begin{enumerate}
\item \emph{Build:} \nolinkurl{lake build MathProofs.SeamB.C042.Solution} completes from a clean module state with exit code~$0$, zero warnings, and 1{,}794 jobs ($5.9\text{ s}$).
\item \emph{Lexical trust-escape scan:} $0$ occurrences of \texttt{sorry}, \texttt{admit}, \texttt{axiom}, \texttt{native\_decide}, \texttt{ofReduceBool}, \texttt{ofReduceNat}, \texttt{implemented\_by}, or \texttt{unsafe} in \texttt{Solution.lean}.
\item \emph{Axiom audit:} \texttt{\#print axioms} confirms that every exported theorem in Table~\ref{tab:lean-map} (all five pinned theorems plus \texttt{a191579\_eq\_a097724}) depends strictly on the standard Lean allowlist \texttt{[propext, Classical.choice, Quot.sound]}.
\item \emph{Independent kernel replay:} Verified via \texttt{lake env leanchecker --fresh --verbose} on module \nolinkurl{MathProofs.SeamB.C042.Solution} (exit code~$0$; full log recorded in the lane verification record).
\end{enumerate}

\begin{thebibliography}{99}
\raggedright

\bibitem{oeis191579}
K.~J.\ Ramsey, \emph{OEIS sequence \seq{A191579}}: Triangular array related to continued fractions of square root of $(N^2 - 1)$ for $N > 1$, apparently containing \seq{A004148} and summing to \seq{A091964}. Submitted June~7, 2011; live revision \#29 (October~13, 2025). \url{https://oeis.org/A191579}.

\bibitem{oeis192062}
K.~J.\ Ramsey (ed.\ O.\ G\'erard), \emph{OEIS sequence \seq{A192062}}: Square array $T(i, j)$ with columns $2j$ being the denominators of continued fraction convergents to $\sqrt{j^2 + 2j}$. Submitted June~21, 2011; live revision \#23 (December~10, 2016). \url{https://oeis.org/A192062}.

\bibitem{oeis097724}
E.\ Deutsch, \emph{OEIS sequence \seq{A097724}}: Triangle read by rows counting left factors of Motzkin paths without peaks, having length $n$ and endpoint height $k$. Includes comments by P.\ Del\'eham (January~26, 2014) and P.\ Bala (June~26, 2025); live revision \#47 (May~30, 2026). \url{https://oeis.org/A097724}.

\bibitem{oeis091964}
E.\ Deutsch, \emph{OEIS sequence \seq{A091964}}: Number of left factors of peakless Motzkin paths of length $n$. Live revision \#68 (May~30, 2026). \url{https://oeis.org/A091964}.

\bibitem{oeis004148}
N.~J.~A.\ Sloane, \emph{OEIS sequence \seq{A004148}}: Generalized Catalan numbers (RNA secondary structure numbers / peakless Motzkin excursions). Live revision \#407 (August~16, 2026). \url{https://oeis.org/A004148}.

\bibitem{oeis089735}
E.\ Deutsch, \emph{OEIS sequence \seq{A089735}}: Self-convolution of \seq{A004148} with itself. Live revision \#27 (November~18, 2026). \url{https://oeis.org/A089735}.

\bibitem{oeis098075}
E.\ Deutsch, \emph{OEIS sequence \seq{A098075}}: Threefold convolution of \seq{A004148} with itself. Live revision \#22 (November~18, 2026). \url{https://oeis.org/A098075}.

\bibitem{stein1979}
P.~R.\ Stein and M.~S.\ Waterman, \emph{On some new sequences generalizing the Catalan and Motzkin numbers}, Discrete Math.\ \textbf{26} (1979), 261--272.

\bibitem{hofacker1998}
I.~L.\ Hofacker, P.\ Schuster, and P.~F.\ Stadler, \emph{Combinatorics of RNA secondary structures}, Discrete Appl.\ Math.\ \textbf{88} (1998), 207--237.

\bibitem{cameron2005}
N.~T.\ Cameron and A.\ Nkwanta, \emph{On some (pseudo) involutions in the Riordan group}, J.\ Integer Seq.\ \textbf{8} (2005), Article 05.3.7.

\bibitem{asinowski2019}
A.\ Asinowski, A.\ Bacher, C.\ Banderier, and B.\ Gittenberger, \emph{Analytic combinatorics of lattice paths with forbidden patterns, the vectorial kernel method, and generating functions for pushdown automata}, Algorithmica \textbf{82} (2020), 386--428 (online 2019).

\bibitem{barry2022}
P.\ Barry, \emph{A three-parameter family of involutions in the Riordan group defined by orthogonal polynomials}, arXiv:2207.02719 (2022). \url{https://arxiv.org/abs/2207.02719}.

\bibitem{prodinger2023}
H.\ Prodinger, \emph{Peakless Motzkin paths of bounded height}, arXiv:2308.03080 (2023). \url{https://arxiv.org/abs/2308.03080}.

\bibitem{prodinger2025}
H.\ Prodinger, \emph{Cornerless, peakless, valleyless Motzkin paths (regular and skew) and applications to bargraphs}, Aequationes Math.\ (2025), arXiv:2501.13645. \url{https://arxiv.org/abs/2501.13645}.

\bibitem{tenn2026}
M.\ Tenn, \emph{Lattice walks, pattern statistics, and Riordan arrays}, arXiv:2607.29612 (2026). \url{https://arxiv.org/abs/2607.29612}.

\end{thebibliography}

\end{document}
